% Transcription — fonds Alexandre Grothendieck, shelfmark 60, batch 3 (pages 41-60).
% Established from the facsimile archives/batches/60/batch-03.pdf (PDF page k =
% archivists' page 40+k, no cover sheet), cut from a copy of 60.pdf fetched
% through the project's relay
% (https://grothendieck-relay.onrender.com/source/60.pdf); tiles in
% archives/tiles/60/p41-p60. Per the skill transcribe-grothendieck. The folder
% is sixty-nine pages in four batches (1-20, 21-40, 41-60, 61-69), transcribed
% in parallel; this is the third. Batch 2 (transcripts/60/batch-02.fr.tex) was
% read first and its notation is kept.
%
% Pass: Opus 5.5 (claude-opus-5-5), 2026-09-26 — first pass, unchecked against the pages by a human.
%
% Pages skipped: 41, 43, 45, 47, 49, 51, 53, 55, 57, 59 — the typed versos of
% the sheets (Bourbaki rédaction n° 392, pages 4-34 to 4-42 and 6-24, 6-25),
% with nothing in his hand: skipped without description by the settled rule.
%
% Pages summarised: none. Pages transcribed: 42, 44, 46, 48, 50, 52, 54, 56,
% 58, 60 — ten rectos in ink. His own pagination: none visible (the pencilled
% numbers bottom left are the archivists'). No numbered formulas; his only
% labelled run is the list of problems a), b), c) on p. 54, which p. 58
% refers back to (« c) est ramené ... au cas de b) »). Boxed formulas are
% kept boxed. Fast drafting: formulas carry, the connective prose of pp. 42,
% 44, 46 and 60 much less so.
%
% What the batch is about. It continues the « Suite de la remarque » of
% batch 2 (the averaging E -> E^I, written E^natural). P. 42: the image of
% R(D/U,k) in R(D,k); cl(E)^natural defined as soon as some open U of finite
% index in I acts unipotently; continuous representations; a struck
% paragraph on D/I = Z. P. 44: the trace formula Tr_{E^natural}(phi) =
% (1/N) sum over the lifts psi_alpha in D/U of phi of Tr_E(psi_alpha), via
% the projector pi = (1/N) sum_{I/U} e_i. P. 46: geometric case: D = D_x,
% I = I_x, D_x/I_x = Z^ (generator f_x); the virtual sheaf E_eta^natural =
% i_!(E_U) + sum_{x in Y} j_{x*}(E_eta^{natural(x)}). Pp. 48-50: its dual,
% D(E_eta^natural) = E_eta^natural-check(1) + sum_x j_{x*}(mu_x(E)), with
% alpha_x(E_eta) = E_eta^{natural(x)} - E_eta(x) and mu_x(E) =
% alpha_x(E_eta)^v - alpha_x(E_eta-check)(1); NB alpha_x is not additive,
% mu_x is; with weight rho, (rho+1) twists. P. 52: the resulting functional
% equation L_{E_eta}(1/qt) = prod lambda_x(E)(t) delta^natural^{-1}
% (-qt)^{chi^natural} L_E(t), q = p^{1+rho}, lambda_x = L_{mu_x}, with a
% bracketed doubt « delta^natural(E_eta)^2 = q^{chi^natural} ?? ». P. 54: the
% same written as L(1/qt) = A(t) L(t), local factors A_x; NB that writing
% lambda_x as a monomial loses additivity; « Pb » a) construct alpha_x in
% the classical cases, b) a general (Shafarevich-)Ogg formula for
% chi^natural, c) a formula for delta^natural (geometric vs arithmetic
% invariants). P. 56: from the dual formula, epsilon^natural(E_eta)^2 as a
% product of local terms, epsilon^natural = ± prod epsilon_x p^{-chi/2}.
% P. 58: c) reduces (up to a sign) to b), i.e. to chi^natural, known via
% Artin representations. P. 60: « cohomologie à l'infini »: R^i_infty f(F)
% = R^i fbar_*(R i_*(F)|Y), the long exact sequence
% R^i f_! -> R^i f_* -> R^i_infty f -> R^{i+1} f_!, and a second triangle
% R j^! R i_! F -> R i_! F -> R i_* F on Xbar.
%
% Notation, as in batch 2: the averaging operation is written \natural (his
% glyph is shaped like a lightning flash); E_\eta^{\natural(x)} for the local
% average at x, \chi^\natural, \delta^\natural, \varepsilon^\natural; \check E
% for the contragredient; D is both the dualising functor and (p. 42-46) the
% decomposition group, as on the page; R without blackboard bold, as in
% batches 1-2 (he writes an outline R on pp. 48 and 60). The letter written
% \chi_x on pp. 52-56 is overwritten and may be \beta_x; one note each.
%
% Readings to check first: the prose of pp. 42 and 44; the Q_l on p. 48
% (mu_x); the \rho read where the ink looks like \varphi on pp. 50-56; the
% name « Chafarevitch » on p. 54; the verbs of the last paragraph of p. 60.

\documentclass[11pt,a4paper]{article}
% Le préambule commun à toutes les transcriptions du fonds Grothendieck.
%
% Il tient en une idée : **l'incertitude se compose, elle ne se résout pas.**
% Une transcription qui lisse un mot douteux a détruit la seule chose pour
% laquelle on la faisait — un lecteur doit pouvoir savoir, à chaque endroit,
% ce qui était sur le feuillet et ce qui a été ajouté. Les six macros qui
% suivent sont donc l'appareil critique, et non de la décoration.
%
% Les mêmes commandes sont reconnues par `scripts/render.mjs`, qui produit la
% vue de lecture du site. Une macro ajoutée ici doit l'être aussi là-bas,
% faute de quoi le rendu échouera bruyamment — ce qui est voulu : un rendu
% silencieusement faux est pire qu'un rendu absent.

\ProvidesPackage{grothendieck}[2026/08/08 Transcriptions du fonds Grothendieck]

\usepackage{amsmath,amssymb,amsthm}
\usepackage{mathrsfs}
\usepackage{stmaryrd}
% Les diagrammes commutatifs sont partout dans ce fonds : c'est la langue
% même de la géométrie algébrique de Grothendieck, et une transcription qui
% les rendrait en prose ne transcrirait rien.
\usepackage{tikz-cd}
% Français par défaut — c'est la langue des feuillets — mais l'anglais est
% chargé aussi : la traduction et le résumé sont des documents anglais, et la
% typographie française (espaces avant les deux-points) y serait une faute.
% Ces documents basculent par \selectlanguage{english} en tête de corps.
\usepackage[english,french]{babel}
\usepackage{fontspec}
\usepackage{geometry}
\geometry{a4paper,margin=25mm,marginparwidth=28mm}
\usepackage{marginnote}
\usepackage{xcolor}
% ulem, not soul. `soul`'s \st and \ul are built on a character-by-character
% scan that XeLaTeX's font machinery breaks: the document fails with an
% undefined control sequence rather than degrading. ulem does the same two jobs
% and is Unicode-engine safe. `normalem` stops it hijacking \emph, which the
% transcriptions use for Grothendieck's own emphasis.
\usepackage[normalem]{ulem}
\usepackage{hyperref}
\hypersetup{colorlinks=true,linkcolor=black,urlcolor=black,pdfborder={0 0 0}}

% --- Métadonnées du lot -----------------------------------------------------
% Reprises telles quelles par le script de rendu, qui les affiche en tête de
% la vue de lecture. Elles fixent ce que le fichier prétend couvrir : sans
% elles, une transcription flotte sans point d'attache dans un fonds de seize
% mille feuillets.

\newcommand{\folder}[1]{\gdef\grothFolder{#1}}
\newcommand{\batch}[1]{\gdef\grothBatch{#1}}
\newcommand{\leaves}[2]{\gdef\grothFirst{#1}\gdef\grothLast{#2}}
\newcommand{\foldertitle}[1]{\gdef\grothFoldertitle{#1}}
\gdef\grothFolder{?}\gdef\grothBatch{?}\gdef\grothFirst{?}\gdef\grothLast{?}\gdef\grothFoldertitle{}

% --- Le résumé --------------------------------------------------------------
%
% Il ouvre la lecture modernisée, sous le titre « Résumé », et il est composé
% en petit. Ce n'est pas un ornement : c'est le seul passage du document qui
% ne soit pas des mathématiques, et un lecteur doit pouvoir le reconnaître
% avant de l'avoir lu — puis le sauter s'il connaît déjà le sujet, ou s'y
% arrêter s'il ne le connaît pas. Le filet qui le suit marque la reprise du
% corps.

\newenvironment{resume}
  {\small\setlength{\parskip}{0.45em}}
  {\par\medskip\hrule\medskip}

% --- Mots-clés ---------------------------------------------------------------
%
% En anglais, en clôture du résumé de la lecture modernisée : le vocabulaire
% moderne sous lequel chercher ce que les feuillets construisent. Le manifeste
% du site (`npm run manifest`) les extrait pour étiqueter la cote entière —
% cette ligne est la source unique des tags, il n'y a pas de fichier à côté.
\newcommand{\keywords}[1]{\par\smallskip\noindent
  {\footnotesize\textsc{Keywords} --- \textit{#1}}\par}

% --- L'appareil critique ----------------------------------------------------

% Le numéro de feuillet, en marge. C'est l'ancre qui permet de régler un
% désaccord sur une formule en regardant *un* feuillet plutôt qu'une chemise
% de six cents. Le site s'en sert aussi pour tourner les pages du fac-similé
% quand on fait défiler la transcription.
\newcommand{\leaf}[1]{%
  \marginnote{\footnotesize\color{gray!70}\textsf{#1}}%
  \label{leaf:#1}\ignorespaces}

% Illisible : on ne devine pas. Un mot inventé qui se lit comme les autres est
% le pire résultat possible.
\newcommand{\ill}{\textcolor{red!60!black}{[\,\ldots\,]}}

% Lecture douteuse : le mot est donné, et signalé comme incertain.
\newcommand{\uncertain}[1]{\uline{#1}}

% Ajout de l'éditeur — une lettre manquante, un mot sous-entendu.
\newcommand{\add}[1]{\textcolor{blue!50!black}{[#1]}}

% Biffé par Grothendieck lui-même. Ce qu'il a rayé fait partie du document :
% une rature dit souvent où la pensée a changé de direction.
\newcommand{\struck}[1]{\sout{#1}}

% Note du transcripteur, sur l'état matériel du feuillet.
\newcommand{\note}[1]{\par\smallskip{\small\color{gray!60!black}\textit{[#1]}}\par\smallskip}

% Note marginale de Grothendieck, distincte du corps du texte.
\newcommand{\marginal}[1]{\par\smallskip\noindent\hspace{1em}%
  {\small\color{gray!60!black}#1}\par\smallskip}

% --- Titre ------------------------------------------------------------------

\AtBeginDocument{%
  \begin{center}
    {\large\bfseries Fonds Alexandre Grothendieck --- cote n\textsuperscript{o}~\grothFolder}\\[2pt]
    {\itshape\grothFoldertitle}\\[4pt]
    {\small feuillets \grothFirst--\grothLast\ (lot \grothBatch)}\\[2pt]
    {\footnotesize\color{gray!60!black}Transcription établie d'après le fac-similé numérisé
     par l'Université de Montpellier.\\
     Les lectures incertaines sont soulignées, les illisibles notées [\,\ldots\,],
     les ajouts entre crochets.}
  \end{center}
  \vspace{1em}}

\endinput


\folder{60}
\batch{3}
\pages{41}{60}
\dating{[à partir de 1962- à partir de 1971]}
\watermark{Édition de démonstration}
\foldertitle{Fonctions L / Équation fonctionnelle des fonctions L (cas géométrique) : notes manuscrites (s.d.)}

\begin{document}

\page{42}
D'autre part, on a
\[
  R(D/U, k) \longrightarrow R(D, k)
\]
et l'image de $R(D/U, k)$ dans $R(D, k)$ \ill{} \uncertain{précisément}
formé\add{e} des classes de $E$ satisfaisant la condition \ill{}. Pour un
tel $E$, on a donc une opération $\natural$ sur la \emph{classe} de $E$.
Donc $\mathrm{cl}(E)^{\natural}$ est défini dès qu'il existe un sous-groupe
\struck{\ill{}} $U$ d'indice fini dans $I$ tel que \struck{$D/U$} $U$ opère
sur $E$ de façon unipotente.
\note{la barre oblique qui suit « unipotente » est de sa main et sépare les
deux alinéas}
Si on prend \ill{} \ill{} de repr.\ \emph{continues} (\ill{} \ill{}
\ill{} corps top.) on \struck{prendra} \add{obtiendra} \ill{} \ill{}
$\mathrm{cl}(E)^{\natural}$ comme \ill{} de repr.\ \emph{continues} de
$D/I$. Compatibilité avec ext.\ des corps de base.
\note{« obtiendra » est écrit au-dessus du mot biffé ; le signe qui précède
$\mathrm{cl}(E)^{\natural}$ ressemble à un $=$ surchargé}

\struck{Supposons que $D/I \simeq \mathbb{Z}$, et donc toute \ill{} \ill{}
de repr.\ \ill{} $D/I$ \ill{} déterminé par sa série form.\
$L_{E^{\natural}}(t)$. Comme \ill{} \ill{}}
\note{ce dernier alinéa, au bas de la page, est barré de trois grands
traits verticaux}

\page{44}
Calculons la \emph{trace} de la représentation \ill{}, on trouve pour
$\varphi \in D/I$
\[
  \operatorname{Tr}_{E^{\natural}}(\varphi) = \sum_i \operatorname{Tr}_{E_i^{\natural}}(\varphi)
\]
\note{après le premier membre, un symbole est noirci ; le $\varphi$ de
« pour $\varphi \in D/I$ » surcharge une autre lettre}
Or si \struck{\ill{}} $U$ opère trivialement sur $E$, \uncertain{opérons}
alors \ill{} \uncertain{modules} sur $D/U$, \ill{} \ill{}
\[
  \ill{}\ E^{\natural} = \operatorname{Im} \pi , \qquad
  \pi = \frac{1}{N} \sum_{i \in I/U} e_i , \qquad N = \operatorname{Card} I/U
\]
\note{sous $\operatorname{Im}\pi$, une lettre biffée et un $L$ indicé d'un
mot, peut-être « régulière » ; l'indice de $\pi$ est biffé}
donc si $g \in I/U \mapsto \varphi$ \ldots\ \ldots
\[
  \operatorname{Tr}_{E^{\natural}}(\varphi) = \operatorname{Tr}_E(\pi g)
  \struck{= \frac{1}{N} \sum \operatorname{Tr}_E(\ill{})}
\]
\note{la suite de la ligne, une somme indicée par $g \mapsto \varphi$, est
effacée de boucles ; à droite, un mot isolé, peut-être « d'où »}
\[
  \boxed{\operatorname{Tr}_{E^{\natural}}(\varphi) = \frac{1}{N}
    \sum_{\substack{\psi_\alpha \in D/U \\ \psi_\alpha \mapsto \varphi \in D/I}}
    \operatorname{Tr}_E \psi_\alpha}
\]
Comme \struck{cette} \add{les deux membres de la} formule \struck{\ill{}}
\uncertain{linéaire} en $E$, \ill{} \ill{} \uncertain{valable} pour
\uncertain{tout} $E$ ; il faut simplement noter que $\operatorname{Tr}_E \psi$
\ill{} défini \struck{\ill{}} pour $\psi \in D/U$, si $U$ opère de façon
unipotente, comme $\operatorname{Tr}_E(\psi')$ pour $\psi' \in D$ se
\uncertain{relevant} en $\psi$ : cette valeur ne \emph{dépend pas} du
\note{l'ajout est écrit au-dessus de la ligne, dans une boucle ; la phrase
continue page 46}

\page{46}
choix de $\psi'$ relevant $\psi$, comme on \uncertain{voit} \ill{} :
\ill{} \ill{} la suite de composition.

Revenons au cas géométrique, on prendra pour $D$ le groupe de
décomposition $D_x$ d'un pt $x \in X$, $I$ le groupe d'inertie, donc
$D_x/I_x \simeq \pi_x \simeq \hat{\mathbb{Z}}$ (générateur $f_x$). On
prendra
\[
  E_\eta(x) = E_\eta^{\natural(x)} \quad \text{comme module sur } \pi_x .
\]
\struck{On posera} Il vient alors, comme \ill{} de \ill{} faisceau virtuel
sur $X$, on posera
\[
  E_\eta^{\natural} = i_{!}(E_U) + \sum_{x \in Y} j_{x*}\bigl(E_\eta^{\natural(x)}\bigr)
\]
où $U$ est un ouvert dense de $X$ sur lequel $E_\eta$ est non ramifié,
$E_U$ est l'image de $E_\eta$ sur $U$, ($i : U \to X$ \add{$Y = X - U$,}
\add{et $j : Y \to X$} désignant les immersions canoniques),
$j_x : \operatorname{Spec} k(x) \to X$.
\note{les deux ajouts sont écrits au-dessus de la ligne, sur des arcs}

\page{48}
Calculons \uncertain{le} $D(E_\eta^{\natural})$ :
\[
  \begin{aligned}
    D(E_\eta^{\natural}) &= D\bigl(i_{!}(E_U)\bigr)
      + \sum_{x \in Y} D\bigl(j_{x*}(E_\eta^{\natural(x)})\bigr) \\
    &= R i_{*}(D E_U) + \sum_{x \in Y} j_{x*}\bigl(D(E_\eta^{\natural(x)})\bigr)
  \end{aligned}
\]
Or la formation de $\natural$ local commute au passage au dual, d'autre
part on a
\[
  D(E_U) = \check E_U(1)[2] \quad \struck{E_U(\rho+1)} , \quad \text{d'où}
\]
\note{le $[2]$ est surchargé et comme encadré ; l'expression qui suit est
barrée de plusieurs traits}
\[
  \begin{aligned}
    D(E_\eta^{\natural}) &= \bigl(R i_{*}(\check E_U)\bigr)(1)
      + \sum_{x \in Y} j_{x*}\bigl(\check E_\eta^{\natural(x)}\bigr) \\
    &= R i_{!}(\check E_U)(1) + \sum_{x \in Y} j_{x*}\Bigl[\check E_\eta^{\natural(x)}
      + j_x^{*} R i_{*}(\check E_U)(1)\Bigr] \\
    &= \check E_\eta^{\natural}(1) + \sum_{x \in Y} j_{x*}\bigl(\mu_x(E)\bigr)
  \end{aligned}
\]
\note{dans les deux premières lignes, devant le $1$ de $(1)$, un signe est
noirci ; au-dessus de $R i_{!}$ un mot biffé, et un $=$ isolé sur la ligne
précédente ; devant la dernière ligne, « Il y a donc » est biffé ; l'indice
de $\mu$ surcharge une autre lettre}
\[
  \mu_x(E) = \check E_\eta^{\natural(x)} \otimes \bigl(\mathbb{Q}_\ell - \mathbb{Q}_\ell(1)\bigr)
    + j_x^{*} R i_{\eta}(\check E_\eta)(1)
\]
\note{devant le premier $\mathbb{Q}_\ell$, des lettres surchargées et
biffées : la lecture $\mathbb{Q}_\ell$ est \uncertain{douteuse}}
\struck{\ill{}} Or par \uncertain{dualité} locale, \struck{\ill{}}
\[
  j_x^{*}\bigl(R i_{\eta}(\check E_\eta)\bigr) = \check E_\eta(x) - \bigl(E_\eta(x)\bigr)^{\vee}(-1)
\]
\note{l'indice de $\check E_\eta(x)$ porte un astérisque noirci}

\page{50}
Posons
\[
  \boxed{\alpha_x(E_\eta) = E_\eta^{\natural(x)} - E_\eta(x)}
\]
Alors
\[
  \boxed{\mu_x(E) = \alpha_x(E_\eta)^{\vee} - \alpha_x(\check E_\eta)(1)}
\]
[\uncertain{donc}
\[
  \mu_x(\check E) = \alpha_x(\check E_\eta)^{\vee} - \alpha_x(E_\eta)(1)
  = -\mu_x(E)^{\vee}(1) \quad \text{i.e.}
\]
\[
  \struck{\mu_x(\check E) = \mu_x}
\]
\[
  \left\lbrace
  \begin{array}{l}
    \mu_x(\check E) = -\mu_x(E)^{\vee}(1) \\
    \mu_x(E) = -\mu_x(\check E)^{\vee}(1)
  \end{array}
  \right.
\]
]
\note{le crochet ouvert devant « donc » se ferme loin à droite, après les
deux formules en accolade}

Donc on trouve
\[
  \boxed{
  \left\lbrace
  \begin{array}{l}
    D(E_\eta^{\natural}) = \check E_\eta^{\natural}(1)
      + \displaystyle\sum_{x \in X^{(0)}} j_{x*}\bigl(\mu_x(E)\bigr) \\[10pt]
    \mu_x(E) = \alpha_x(E_\eta)^{\vee} - \alpha_x(\check E_\eta)(1) \qquad
    \alpha_x(E_\eta) = E_\eta^{\natural(x)} - E_\eta(x)
  \end{array}
  \right.}
\]

\emph{NB}. $\alpha_x(E_\eta)$ n'est \emph{pas} additif en $E_\eta$,
\emph{mais} $\mu_x(E_\eta)$ \emph{l'est}.
\note{« mais » et « l'est » sont soulignés deux fois}

Si $\check E_\eta = E_\eta(\rho)$, on trouve
\note{ici et dans l'encadré qui suit, la lettre ressemble à un $\varphi$ ;
on lit le $\rho$ des pages précédentes}
\[
  \boxed{
  \begin{array}{l}
    D(E_\eta^{\natural}) = E_\eta^{\natural}(\rho+1)
      + \displaystyle\sum_{x \in X^{(0)}} j_{x*}\bigl(\mu_x(E)\bigr) \\[10pt]
    \mu_x(E) = \alpha_x(E_\eta)^{\vee} - \alpha_x(E_\eta)(\rho+1)
  \end{array}}
\]

\page{52}
L'équation fonctionnelle devient alors
\[
  L_{\check E_\eta}(p^{-1} t) \prod_x \lambda_x(E)
  = (-t)^{-\chi^{\natural}(E_\eta)}\, \delta^{\natural}(E_\eta)\, L_{E_\eta}(t^{-1})
\]
\marginal{à droite : i.e., remplaçant $t$ par $pt$}
\note{le $\lambda_x$ surcharge une autre lettre}
\[
  \left\lbrace
  \begin{array}{l}
    L_{E_\eta}\Bigl(\dfrac{1}{pt}\Bigr) = \prod_x \lambda_x(E)\,
      (-pt)^{\chi^{\natural}(E_\eta)}\, \delta^{\natural}(E_\eta)^{-1}\, L_{\check E_\eta}(t) \\[10pt]
    \chi^{\natural}(E_\eta) = \chi(E_\eta^{\natural}) \qquad \text{à déterminer} \ldots \\[4pt]
    \delta^{\natural}(E_\eta) = \delta(E_\eta^{\natural}) \qquad \text{à déterminer} \ldots \\[4pt]
    \lambda_x(E) = L_{\mu_x(E)}(t) = L_{\alpha_x(E_\eta)^{\vee}}(t) / L_{\alpha_x(\check E_\eta)}(p^{-1} t)
      = (-t)^{\chi_x(E_\eta)}\, \delta_x(E_\eta)\,
      \dfrac{L_{\alpha_x(E_\eta)}(t^{-1})}{L_{\alpha_x(\check E_\eta)}(t)}
  \end{array}
  \right.
\]
\note{au premier membre, le signe de contragrédiente sur $E_\eta$ paraît
biffé d'un point ; une accolade réunit $(-pt)^{\chi^{\natural}(E_\eta)}$ et
$\delta^{\natural}(E_\eta)^{-1}$ ; dans la dernière ligne, l'exposant et
l'argument de $\delta_x$ sont surchargés et noircis}

Supposons $E$ de poids $\rho$, donc
\[
  \check E \simeq E(\rho) , \qquad L_{\check E}(t) = L_E(p^{-\rho} t) ,
\]
d'où il vient, en remplaçant $t$ par $p^{\rho} t$
\[
  \boxed{L_{E_\eta}\Bigl(\frac{1}{qt}\Bigr) = \Bigl(\prod_x \lambda_x(E)(t)\Bigr)\,
    \delta^{\natural}(E_\eta)^{-1}\, (-qt)^{\chi^{\natural}(E_\eta)}\, L_E(t)}
\]
\note{devant $q$, dans $(-qt)$, une lettre est biffée}
\[
  \Bigl[\delta^{\natural}(E_\eta)^{-1} q^{\chi^{\natural}(E_\eta)}
    \overset{?}{=} \struck{\varepsilon(E)}\ \delta^{\natural}(E_\eta)
    \quad \text{i.e.} \quad
    \delta^{\natural}(E_\eta)^2 = q^{\chi^{\natural}(E_\eta)} \ ??\Bigr]
    \longleftarrow ??
\]
\note{le point d'interrogation au-dessus du $=$ est entouré ; devant $q$, un
symbole est noirci}
\marginal{à gauche, biffés : $q = p^{\ill{}}$ et $\lambda_x$}
\[
  \left\lbrace
  \begin{array}{l}
    q = p^{1+\rho} \\[4pt]
    \lambda_x(E_\eta) = L_{\mu_x(E_\eta)}(t)
      = (-t)^{-\chi_x(E_\eta)}\, \delta_x(E_\eta)\,
      \dfrac{L_{\alpha_x(E_\eta)}(t^{-1})}{L_{\alpha_x(E_\eta)}(p^{-\rho} t)} \\[10pt]
    \chi_x(E_\eta) = \chi\bigl(\alpha_x(E_\eta)\bigr) \\[4pt]
    \delta_x(E_\eta) = \delta\bigl(\alpha_x(E_\eta)\bigr)
  \end{array}
  \right.
\]
\marginal{à droite de $\lambda_x$ :
$\mu_x(E_\eta) = \alpha_x(E_\eta)^{\vee} - \alpha_x(E_\eta)(\rho+1)$}
\note{un premier membre de droite de $\lambda_x(E_\eta)$ est lourdement
biffé ; la lettre notée ici $\chi_x$, dans l'exposant et dans la définition,
est surchargée et pourrait être un $\beta_x$}

\page{54}
\[
  L_{E_\eta}\Bigl(\frac{1}{qt}\Bigr) = A(t)\, L_{E_\eta}(t)
\]
\[
  A(t) = \prod_x \frac{L_{\alpha_x(E_\eta)}(t^{-1})}{L_{\alpha_x(E_\eta)}(p^{-\rho} t)}\;
    (-t)^{\chi^{\natural}(E_\eta) - \sum_x \chi_x(E_\eta)}\;
    \frac{q^{\chi^{\natural}(E_\eta)}}{\delta^{\natural}(E_\eta) \prod_x \delta_x(E_\eta)^{-1}}
\]
\note{la fraction sous le produit est entourée ; un trait la relie à
« $= A_x(t^{-1})^{-1}$ », écrit au-dessous}
\[
  L_{\alpha_x(E_\eta)}\Bigl(\frac{1}{p^{\rho} t}\Bigr) = A_x(t)\, L_{\alpha_x(E_\eta)}(t)
\]
\note{au-dessous, un début de ligne, $L_{\alpha_x(E_\eta)} \ldots
L_{\alpha_x(E_\eta)}($, est biffé}

N.B. Ayant écrit $\lambda_x(E_\eta)(t)$ comme un \uncertain{monôme}, on en
perd le \emph{caractère additif}, il vaut sans doute mieux ne pas faire ce
\uncertain{changement} \ldots
\note{l'argument de $\lambda_x$ surcharge une première écriture}

\emph{Pb} a) Dans les cas géométriques classiques, construire les
$\alpha_x(E_\eta)$ qui déterminent les $\mu_x$, donc les
$\lambda_x(E_\eta)(t)$. Déterminer en particulier $\chi_x(E_\eta)$,
$\delta_x(E_\eta)$, les valeurs absolues des val.\ propres de Frobenius.

b) Donner une formule générale (\uncertain{Chafarevitch}--Ogg) pour
$\chi^{\natural}(E_\eta)$

c) Donner une formule pour $\delta^{\natural}(E_\eta)$
\marginal{des flèches relient $\chi_x(E_\eta)$ et $\chi^{\natural}(E_\eta)$
à « invariant géom. », $\delta_x(E_\eta)$ et $\delta^{\natural}(E_\eta)$ à
« invariant arithm. », ces deux mots soulignés deux fois}

\page{56}
Sur la formule virtuelle donnant $D(E_\eta^{\natural})$, on lit
\[
  \delta^{\natural}(E_\eta)^{-1} = \delta^{\natural}(E_\eta)\,
    p^{-\chi^{\natural}(E)(\rho+1)} \prod_x \delta_x\bigl(\mu_x(E)\bigr)
\]
Or
\[
  \begin{aligned}
    \delta\bigl(\mu_x(E)\bigr) &= \delta\bigl(\alpha_x(E_\eta)\bigr)^{-1}
      \Bigl[\delta\bigl(\alpha_x(E_\eta)\bigr)\, p^{-\chi_x(E_\eta)(\rho+1)}\Bigr]^{-1} \\
    &= \delta_x(E_\eta)^{-2}\, q^{+\chi_x(E_\eta)}
  \end{aligned}
\]
d'où
\[
  \Bigl\lbrace \prod_x \frac{\delta_x(E_\eta)^2}{q^{\chi_x(E_\eta)}} \Bigr\rbrace
  = \delta^{\natural}(E_\eta)^2\, q^{\sum \chi_x(E_\eta)}
\]
\note{l'accolade ouvrante est surchargée}
\[
  \left\lbrace
  \begin{array}{l}
    \dfrac{\delta^{\natural}(E_\eta)^2}{q^{\chi^{\natural}(E)}}
      = \displaystyle\prod_x \frac{\delta_x(E_\eta)^2}{q^{\chi_x(E_\eta)}} \\[14pt]
    \qquad \parallel \\[2pt]
    \varepsilon^{\natural}(E_\eta)^2
      = \displaystyle\prod_x \frac{\delta_x(E_\eta)^2}{q^{\chi_x(E_\eta)}}
  \end{array}
  \right.
\]
\note{après $\delta^{\natural}(E_\eta)^2$, un premier « $= q^{\chi^{\natural}(E)}$ »
est biffé d'un trait oblique et remplacé par le dénominateur ; à droite de
la première ligne, un mot, peut-être « on aura », ne se lit pas sûrement}
\[
  \varepsilon^{\natural}(E_\eta)
  = \pm \prod_x \frac{\delta_x(E_\eta)}{q^{\chi_x(E_\eta)/2}}
  = \Bigl(\prod_x \frac{\delta_x(E_\eta)}{p^{\rho\chi_x(E_\eta)/2}}\Bigr) \frac{1}{p^{\chi/2}}
\]
\[
  \varepsilon^{\natural}(E_\eta) = \frac{\delta^{\natural}(E_\eta)}{q^{\chi^{\natural}(E_\eta)/2}}
  \qquad\qquad
  \chi = \sum_x \chi_x(E_\eta)
\]
\note{le signe $\parallel$ vertical, sous $\varepsilon^{\natural}(E_\eta)$,
renvoie à sa définition ; deux flèches mènent de $\chi = \sum_x \chi_x(E_\eta)$
aux exposants $\chi/2$}
\[
  \varepsilon^{\natural}(E_\eta) = \pm \Bigl(\prod_x \varepsilon_x(E_\eta)\Bigr) p^{-\chi/2}
\]
\note{la lettre notée $\chi_x$ sur cette page, comme page 52, est
surchargée à plusieurs endroits et pourrait être un $\beta_x$}

\page{58}
On voit donc que c) est ramené (\emph{moyennant un signe}) au cas de b),
donc à déterminer $\chi^{\natural}(E_\eta)$. Or cette détermination est
faite, en utilisant les représentations d'Artin \ldots
\note{« moyennant un signe » est souligné ; le reste de la page est
blanc}

\page{60}
Cohomologie à l'infini :
\note{la fin de cette ligne porte quelques points de suspension, à droite}

$R^{i}_{\infty} f(F)$ pour $f$ compactifiable en $\bar f$.
On \uncertain{envisage} la suite exacte

\begin{tikzcd}
R i_{!}(F) \arrow[r] & R i_{*}(F) \arrow[r] & R i_{*}(F)|Y \arrow[ll, bend right=20]
\end{tikzcd}

\note{la flèche courbe, du troisième terme vers le premier, est le
morphisme de degré $+1$ du triangle ; l'indice du premier terme est
surchargé ($!$ sur $*$)}
d'où en \uncertain{appliquant} $R \bar f_{*}$ une suite exacte
\[
  \boxed{\cdots \to R^{i} f_{!}(F) \to R^{i} f_{*}(F) \to R^{i}_{\infty} f(F)
    \to R^{i+1} f_{!}(F) \to \cdots}
\]
en posant
\[
  R^{i}_{\infty} f(F) = R^{i} \bar f_{*}\bigl(R i_{*}(F)|Y\bigr)
\]
On peut aussi considérer la suite exacte de \uncertain{dualisation}
\note{ce dernier mot est écrit au-dessus de la ligne, sous « On peut »}
\uncertain{pour} $R i_{!}(F)$ sur $\bar X$, et l'ouvert $X$ et
$Y = \bar X - X$, \struck{\ill{}} \uncertain{\ill{}} \ill{} :
\note{sous « $R i_{!}(F)$ », un $R i_{*}(F)$ est biffé ; le $X$ de « l'ouvert
$X$ » surcharge une autre lettre ; la fin de la ligne est surchargée et se
lit mal}

\begin{tikzcd}
R j^{!}\bigl(R i_{!}(F)\bigr) \arrow[r] & R i_{!}(F) \arrow[r] & R i_{*}(F) \arrow[ll, bend right=20]
\end{tikzcd}

\note{au-dessous, une ligne commencée
« $R^{i} \ldots R^{i} f_{!}(F) \to R^{i} f_{*}(F) \to$ » est biffée d'un
long trait}
et on \uncertain{identifie} d'une \uncertain{manière} \uncertain{précise},
$R i_{*}(F)|Y$ à \uncertain{décalé} près \uncertain{par}
$R j^{!} R i_{!}(F)$ \ldots
\note{la page s'arrête sur un tiret}

\end{document}
